% AI-assisted manuscript source; responsible author: Carptopus.
\documentclass[11pt]{article}
\usepackage[a4paper,margin=29mm]{geometry}
\usepackage[T1]{fontenc}
\usepackage[utf8]{inputenc}
\usepackage{lmodern}
\usepackage{microtype}
\usepackage{amsmath,amssymb,amsthm,mathtools}
\usepackage{enumitem}
\usepackage{xcolor}
\usepackage[colorlinks=true,linkcolor=blue!55!black,citecolor=blue!55!black,urlcolor=blue!65!black]{hyperref}
\usepackage{fancyhdr}

\newtheorem{theorem}{Theorem}[section]
\newtheorem{lemma}[theorem]{Lemma}
\newtheorem{proposition}[theorem]{Proposition}
\newtheorem{corollary}[theorem]{Corollary}
\theoremstyle{remark}
\newtheorem{remark}[theorem]{Remark}
\newtheorem*{maintheorem}{Theorem A}
\numberwithin{equation}{section}

\pagestyle{fancy}
\fancyhf{}
\fancyhead[L]{Carptopus}
\fancyhead[R]{Comb-bounded crowns}
\fancyfoot[C]{\thepage}
\setlength{\headheight}{14pt}
\setlength{\parindent}{0pt}
\setlength{\parskip}{0.55em}
\setlength{\emergencystretch}{4em}
\setlist{nosep,leftmargin=*}

\title{Comb-bounded crowns and bi-Esakia representability\\of well-ordered rooted trees}
\author{Carptopus\\\href{mailto:carptopus@163.com}{carptopus@163.com}}
\date{21 August 2026}

\hypersetup{
  pdftitle={Comb-bounded crowns and bi-Esakia representability of well-ordered rooted trees},
  pdfauthor={Carptopus},
  pdfsubject={A finite-anchor synchronization theorem and a bi-Esakia representability criterion for comb-bounded crowns},
  pdfkeywords={bi-Esakia space, bi-Heyting algebra, well-ordered tree, bi-p-morphism, inverse limit, better partial order, finite comb}
}

\begin{document}
\maketitle

\begin{center}
\small\itshape
Public Beta v0.1. The corrected proof and the complete manuscript have passed
the repository's internal writing-admission, manuscript, and artifact audits.
External mathematical review and journal peer review are pending.
\end{center}
\vspace{0.8em}

\begin{abstract}
We give a sufficient condition for an abstract well-ordered rooted tree to
admit a bi-Esakia topology. Let $X$ have root $r$, and for every cover $a$ of
$r$ write $Q_a=\uparrow a$ for the corresponding first-level crown. We prove
that $X$ is bi-Esakia representable whenever all $Q_a$ are finite and their
order duals uniformly omit one fixed finite co-tree comb.

The proof separates topology from finite order theory. First, we prove a
synchronization theorem for arbitrarily many finite posets amalgamated over a
common finite downset. If the petals fall into finitely many classes and every
petal admits a stem-exact surjective bi-p-morphism onto an actual anchor petal,
then the amalgam carries a bi-Esakia topology. The topology is a finite sum of
one-point compactifications of anchor fibers and also arises as the
inverse-limit topology of the finite-petal quotients. Second, a
better-partial-order theorem of Martins supplies finitely many actual minimal
anchor types for every uniformly comb-bounded family of finite crowns. The
resulting class allows arbitrary root width, no uniform bound on crown size,
and no uniform depth bound. We also explain why a known infinite antichain of
finite co-trees obstructs the finite-anchor method without providing a
nonrepresentability theorem.
\end{abstract}

\section{Introduction}

Esakia duality represents Heyting algebras by ordered Stone spaces in which the
down-closure of every clopen set is clopen. Bi-Heyting algebras require the same
topology to support both directions: both down-closures and up-closures of
clopen sets must remain clopen. Consequently, showing separately that a poset
and its order dual are Esakia representable does not by itself produce a
bi-Esakia topology on the original poset.

Fornasiere and Moraschini characterized the Esakia-representable well-ordered
forests satisfying the appropriate chain-supremum condition \cite{ref4}. Their
recursive topology is intrinsically one-sided, and the corresponding bi-Esakia
representability problem remains subtler. Tasiou proved inverse-limit closure
for bi-Esakia spaces in her master's thesis \cite[Theorem 4.2.6]{ref1},
synchronized arbitrary disjoint copies of a fixed finite connected poset
\cite[Theorem 4.3.4]{ref1}\cite[Theorem 4.6]{ref2}, and represented the finite
stages of a countable inverse system as the zigzag image-finite part of a
larger bi-Esakia space
\cite[Theorem 4.3.8]{ref1}\cite[Theorem 4.3]{ref2}. The present question is
different: the finite rooted pieces need not be isomorphic, and they meet at a
common root, or more generally at a common finite stem.

Our first result isolates the required topological construction. Suppose finite
petals $P_i$ share a finite downset $K$. If they can be partitioned into
finitely many groups so that every petal in a group maps by a surjective
bi-p-morphism onto an actual anchor petal, with the map fixing $K$ and taking no
external point into $K$, then the amalgam is bi-Esakia representable. The
topology synchronizes all fibers lying over the same anchor point. This extends
the point-synchronization idea for isomorphic copies to finitely many
nonisomorphic quotient types and a common stem.

The second result applies this construction to trees. Martins proved that
finite co-trees omitting a fixed finite comb form a better partial order under
the surjective bi-p-image relation \cite[Proposition 4.10]{ref5}. Thus, among
the finite crowns that actually occur above a root, there are only finitely many
minimal quotient types. Choosing actual crowns as representatives gives
exactly the anchors required by the synchronization theorem.

\begin{maintheorem}
Let $X$ be a well-ordered rooted tree with root $r$. For each
$a\in\operatorname{Cov}_X(r)$, put $Q_a=\uparrow_Xa$. Assume that
\begin{enumerate}
\item every $Q_a$ is finite; and
\item for some fixed $n\geq1$, the finite tree-comb
  $H_n=C_n^{\mathrm{op}}$ does not order-embed into any $Q_a$, where $C_n$ is
  the finite co-tree $n$-comb of \cite{ref5}.
\end{enumerate}
Then the underlying poset $X$ admits a bi-Esakia topology.
\end{maintheorem}

The well-ordered hypothesis is essential to the statement as written. Indeed,
let
\[
X=\{0\}\cup(\mathbb Q\cap(0,1))
\]
with the usual order. It is a rooted tree whose root has no cover, so the crown
assumptions are vacuous. It has no gaps and is not even Priestley representable
by \cite[Proposition 2.5]{ref4}. The first frozen version of Theorem~A omitted
the well-ordered hypothesis and was therefore false. Section~\ref{sec:crowns}
uses well-ordering to prove that the first-level crowns exhaust all nonroot
points.

The conclusion covers trees of arbitrary root width. The crown sizes need not
have a common bound, and their depths may be unbounded. The uniform restriction
is order-theoretic: the family omits one fixed comb.

\section{Preliminaries}

All posets are ordered by $\leq$. For $x$ in a poset $P$, write
$\uparrow_Px$ and $\downarrow_Px$ for its principal up- and downsets. A tree is
a poset with a least element in which every principal downset is a chain. It is
well-ordered if every principal downset is a well-ordered chain.

A map $f:P\to Q$ is a \emph{surjective bi-p-morphism} if it is surjective and
order preserving and, for every $x\in P$,
\begin{equation}
f[\uparrow_Px]=\uparrow_Qf(x),
\qquad
f[\downarrow_Px]=\downarrow_Qf(x).
\label{eq:bip}
\end{equation}
Equivalently, it satisfies both the bounded and co-bounded back conditions.
These conditions imply, for every $D\subseteq Q$,
\begin{equation}
\downarrow_P f^{-1}(D)=f^{-1}(\downarrow_QD),
\qquad
\uparrow_P f^{-1}(D)=f^{-1}(\uparrow_QD).
\label{eq:inverse-cones}
\end{equation}

A bi-Esakia space is an ordered Stone space whose clopen sets have clopen down-
and up-closures. We use the equivalent criterion from
\cite[Theorem 2.12]{ref3}: besides the clopen-cone conditions, it is enough in
the present Stone setting to verify that every principal upset is closed.

For each positive integer $n$, let $C_n$ be the finite co-tree with points
\[
\{c_1,\ldots,c_n,c'_1,\ldots,c'_n\},
\]
whose order is generated by
\begin{equation}
c_1<c_2<\cdots<c_n,
\qquad
c'_i<c_i\quad(1\leq i\leq n).
\label{eq:comb}
\end{equation}
This is the $n$-comb of \cite[Figure 6]{ref5}, and we put
$H_n=C_n^{\mathrm{op}}$. We define the height of a finite poset to be the
maximum number of points in a chain. Thus $C_1=H_1^{\mathrm{op}}$ is a
two-point chain, $H_n$ has height $n+1$, and $H_2$ is branched and contains a
three-point chain.

For finite co-trees, write $M\preceq_p Q$ when there exists a surjective
bi-p-morphism $Q\to M$. Martins \cite{ref5} proves that, for each fixed $n$,
the class
\begin{equation}
\mathcal T_n=
\{Y:Y\text{ is a finite co-tree and }C_n
  \text{ does not order-embed into }Y\}
\label{eq:Tn}
\end{equation}
is a better partial order under $\preceq_p$. In particular, it is a well
partial order.

\section{Synchronization over a common finite stem}

We first prove the topological part independently of trees.

\begin{theorem}[finite-anchor synchronization]\label{thm:synchronization}
Let $K$ be a finite poset and let $G$ be finite. For each $g\in G$, let $I_g$
be an arbitrary index set with a distinguished anchor $0_g\in I_g$. Suppose
that finite posets $P_{g,i}$ satisfy the following conditions and that their
amalgam is nonempty.
\begin{enumerate}
\item Every $P_{g,i}$ contains the same copy of $K$, and $K$ is a downset of
  $P_{g,i}$.
\item Distinct petals meet exactly in $K$, and external points belonging to
  distinct petals are incomparable.
\item For every $i\in I_g$ there is a surjective bi-p-morphism
  \[
  f_{g,i}:P_{g,i}\longrightarrow P_{g,0_g}
  \]
  which is the identity on $K$ and satisfies $f_{g,i}^{-1}(K)=K$. For the
  anchor, take the identity map.
\end{enumerate}
Let $X$ be the amalgam of all petals over $K$. Then $X$ admits a bi-Esakia
topology.
\end{theorem}

\begin{proof}
For $g\in G$ and $a\in P_{g,0_g}\setminus K$, define the synchronized fiber
\begin{equation}
E_{g,a}=\{a\}\cup
\bigcup_{i\in I_g\setminus\{0_g\}}f_{g,i}^{-1}(a).
\label{eq:fiber}
\end{equation}
Give $E_{g,a}$ the one-point compactification topology in which every point
other than $a$ is isolated and $a$ is the unique possible nonisolated point.
Give $K$ the discrete topology. Since $G$, $K$, and every anchor petal are
finite, $X$ is the finite topological sum
\begin{equation}
X=K\mathbin{\dot\cup}
\mathop{\dot\bigcup}_{g\in G}
\mathop{\dot\bigcup}_{a\in P_{g,0_g}\setminus K}E_{g,a},
\label{eq:sum}
\end{equation}
and is therefore a compact Stone space.

A set $U\subseteq X$ is clopen precisely when, for each synchronized fiber,
its trace is finite if the anchor is absent and cofinite if the anchor is
present. Because the number of anchors is finite and every individual fiber
$f_{g,i}^{-1}(a)$ is finite, there is for each $g$ a finite exceptional set
$J_g\subseteq I_g$ such that
\begin{equation}
U\cap P_{g,i}
=f_{g,i}^{-1}(U\cap P_{g,0_g})
\qquad(i\notin J_g).
\label{eq:stable}
\end{equation}
Conversely, this eventual inverse-image condition characterizes clopen
subsets.

We show that the down-closure of $U$ has the same form. Fix a nonexceptional
petal and abbreviate
\[
U_i=U\cap P_{g,i},\qquad U_0=U\cap P_{g,0_g},
\qquad
T=(\downarrow_XU)\cap K.
\]
No external point is comparable across distinct petals, $K$ is a common
downset, and $f_{g,i}^{-1}(T)=T$. Hence
Equation~\eqref{eq:inverse-cones} gives
\begin{align}
(\downarrow_XU)\cap P_{g,i}
&=\downarrow_{P_{g,i}}U_i\cup T \notag\\
&=f_{g,i}^{-1}
  (\downarrow_{P_{g,0_g}}U_0\cup T) \notag\\
&=f_{g,i}^{-1}
  ((\downarrow_XU)\cap P_{g,0_g}).
\label{eq:down-stable}
\end{align}
The up-closure cannot move from an external point into another petal, while
the common influence of $U\cap K$ is already present in every inverse image.
Thus
\begin{align}
(\uparrow_XU)\cap P_{g,i}
&=\uparrow_{P_{g,i}}U_i \notag\\
&=f_{g,i}^{-1}(\uparrow_{P_{g,0_g}}U_0) \notag\\
&=f_{g,i}^{-1}((\uparrow_XU)\cap P_{g,0_g}).
\label{eq:up-stable}
\end{align}
Only the finitely many exceptional finite petals and the finite set $K$ can
deviate from \eqref{eq:down-stable} or \eqref{eq:up-stable}. Therefore
$\downarrow_XU$ and $\uparrow_XU$ are clopen.

It remains to check principal upsets. If $x\in K$, stem-exactness and the
co-bounded back condition imply that each synchronized fiber lies either
wholly inside $\uparrow_Xx$ or wholly outside it. Hence $\uparrow_Xx$ is
clopen. If $x\notin K$, then $\uparrow_Xx$ lies inside one finite petal and is
finite; it is closed because the Stone space is $T_1$. The criterion from
\cite[Theorem 2.12]{ref3} now shows that $X$ is bi-Esakia.
\end{proof}

\begin{proposition}[inverse-limit realization]\label{prop:inverse-limit}
The topology in Theorem~\ref{thm:synchronization} is the inverse-limit topology
of the finite amalgams obtained by retaining finitely many petals and
collapsing every unretained petal onto its anchor.
\end{proposition}

\begin{proof}
Let
\begin{equation}
\mathcal D=
\prod_{g\in G}
\{F_g\in[I_g]^{<\omega}:0_g\in F_g\},
\label{eq:directed-set}
\end{equation}
ordered by coordinatewise inclusion. It is directed because $G$ is finite and
coordinatewise finite unions give common upper bounds. For $F\in\mathcal D$,
let $X_F$ retain exactly the petals indexed by $F_g$ in each group, and give
this finite poset its discrete bi-Esakia topology. If $F\leq E$, define
\begin{equation}
\pi_{EF}:X_E\longrightarrow X_F
\label{eq:bonding}
\end{equation}
to be the identity on $K$ and on petals already retained by $F$, and to use
$f_{g,i}$ on every petal indexed by $E_g\setminus F_g$. The three possible
locations of a point---stem, retained petal, or collapsed petal---show directly
that both back conditions hold. The definitions also give
\begin{equation}
\pi_{FF}=\operatorname{id}_{X_F},
\qquad
\pi_{LF}=\pi_{EF}\circ\pi_{LE}
\quad(F\leq E\leq L).
\label{eq:coherence}
\end{equation}
Thus $(X_F,\pi_{EF})_{F\in\mathcal D}$ is an inverse system of finite posets and
surjective bi-p-morphisms.

Define the canonical quotient $q_F:X\to X_F$ by retaining every $F_g$-petal
and applying $f_{g,i}$ to every other petal. Then
$q_F=\pi_{EF}\circ q_E$ whenever $F\leq E$, so
\begin{equation}
\eta:X\longrightarrow\varprojlim_{F\in\mathcal D}X_F,
\qquad
\eta(x)=(q_F(x))_F,
\label{eq:eta}
\end{equation}
is well defined.

Let $(x_F)_F$ be compatible, and inspect the least stage containing only the
anchors. If its coordinate belongs to $K$, stem-exactness forces every higher
coordinate, and hence every coordinate, to be the same stem point. Otherwise
that coordinate is an external anchor point. Either every coordinate remains
that anchor, or at some stage $F$ a retained nonanchor petal contains an
external point $y$. Above $F$, compatibility and identity on retained petals
lock every coordinate to $y$. For an incomparable stage $H$, a common upper
bound $L\geq F,H$ uniquely determines the $H$-coordinate by projection. Two
distinct petals cannot both be selected, since a stage retaining both would
have to project to two different external points. Thus $\eta$ is surjective.
It is also injective. For $x\ne y$, choose a finite stage retaining the
nonanchor petals to which the external points among $x,y$ belong; there are at
most two such petals. Identity on retained petals separates two external
points, while stem-exactness separates a point of $K$ from every distinct
point. Hence $q_F(x)\ne q_F(y)$.

For topology, $q_F^{-1}(\{k\})=\{k\}$ for $k\in K$, by stem-exactness. The
preimage of a nonanchor retained point is an isolated singleton. The preimage
of an anchor point $a$ contains $a$ and the full $a$-fiber in every unretained
petal, while it differs only within finitely many retained finite fibers. Thus
every $q_F$ is continuous for the synchronized topology, so that topology is
at least as fine as the cylinder topology.

Conversely, let $E_{g,a}\setminus R$ be a basic synchronized neighborhood of
an anchor, where $R$ is finite and does not contain $a$. Choose $F$ retaining
every petal that contains a point of $R$, and put
\begin{equation}
V=\{a\}\cup\bigl((E_{g,a}\cap X_F)\setminus R\bigr).
\label{eq:V}
\end{equation}
Since $X_F$ is discrete, $V$ is open, and direct inspection gives
\begin{equation}
q_F^{-1}(V)=E_{g,a}\setminus R.
\label{eq:cylinder}
\end{equation}
The isolated stem and nonanchor points were already realized as singleton
cylinders. These sets form a base for the finite topological sum
\eqref{eq:sum}, so the synchronized topology is also no finer than the cylinder
topology. Hence the two topologies agree. Since $X$ is compact and the inverse
limit of finite discrete spaces is Hausdorff, the continuous bijection $\eta$
is a homeomorphism.
\end{proof}

\section{Comb-bounded crowns}\label{sec:crowns}

We now prove Theorem~A.

\begin{lemma}[root-crown decomposition]\label{lem:decomposition}
Let $X$ be a well-ordered rooted tree with root $r$. Then
\begin{equation}
X=\{r\}\mathbin{\dot\cup}
\mathop{\dot\bigcup}_{a\in\operatorname{Cov}_X(r)}\uparrow_Xa,
\label{eq:decomposition}
\end{equation}
and external points from distinct crowns are incomparable.
\end{lemma}

\begin{proof}
For $x>r$, the nonempty interval $(r,x]$ is a subset of the well-ordered chain
$\downarrow x$ and therefore has a least element $a$. No point lies strictly
between $r$ and $a$, so $a$ covers $r$, and $x\in\uparrow a$. If two different
covers $a,b$ both lie below $x$, the chain $\downarrow x$ makes them
comparable, contradicting that both cover $r$. The same argument rules out
comparability between external points in different crowns.
\end{proof}

\begin{theorem}[comb-bounded crown theorem]\label{thm:crowns}
Under the hypotheses of Theorem~A, $X$ is bi-Esakia representable.
\end{theorem}

\begin{proof}
Let $A=\operatorname{Cov}_X(r)$ and $Q_a=\uparrow_Xa$. By
Lemma~\ref{lem:decomposition}, the sets $P_a=\{r\}\cup Q_a$ are finite petals
amalgamated over the common downset $K=\{r\}$.

Every dual crown $Q_a^{\mathrm{op}}$ belongs to $\mathcal T_n$. Let $S$ be the
set of isomorphism types that actually occur among these dual crowns. For each
$s\in S$, well-foundedness gives a minimal element of the nonempty set
$S\cap\downarrow s$. It is globally minimal in $S$, since any smaller element
of $S$ would also lie in $S\cap\downarrow s$. The global minimal elements of
$S$ form an antichain and are therefore finite. Choose, from the original tree,
one actual crown index $a_j$ for each minimal isomorphism type and put
\begin{equation}
M_j=Q_{a_j}^{\mathrm{op}}
\qquad(j=1,\ldots,t).
\label{eq:anchors}
\end{equation}
Hence, for every $a\in A$, there is a $j=j(a)$ and a surjective bi-p-morphism
\begin{equation}
g_a:Q_a^{\mathrm{op}}\longrightarrow M_j.
\label{eq:dual-map}
\end{equation}
After transporting along an isomorphism if necessary, the target is the chosen
actual representative. Taking order duals preserves both back conditions, so
\begin{equation}
g_a^{\mathrm{op}}:Q_a\longrightarrow Q_{a_j}
\label{eq:primal-map}
\end{equation}
is again a surjective bi-p-morphism.

Extend \eqref{eq:primal-map} to $f_a:P_a\to P_{a_j}$ by setting $f_a(r)=r$.
This map is root-exact. At the root, the up-back condition follows from
surjectivity of \eqref{eq:primal-map}, while the only lower point is the root
itself. Inside a crown, both back conditions come from
\eqref{eq:primal-map}, except that a lower target equal to the common root is
lifted by the root. Therefore $f_a$ is a surjective bi-p-morphism and
$f_a^{-1}(\{r\})=\{r\}$.

There are only finitely many anchors $P_{a_j}$.
Theorem~\ref{thm:synchronization} applies and supplies a bi-Esakia topology on
$X$.
\end{proof}

\section{Consequences}

\begin{corollary}[three-level trees]
Every well-ordered rooted tree of height at most three in which each first-level
point has finitely many immediate successors is bi-Esakia representable, with
no restriction on the cardinality of the first level.
\end{corollary}

\begin{proof}
Each crown is a finite star and therefore has height at most two. Since $H_2$
has height three, no crown contains $H_2$. Apply Theorem~\ref{thm:crowns}.
\end{proof}

\begin{corollary}[uniform finite height]
Let $X$ be a well-ordered rooted tree whose first-level crowns are finite. If
the heights, measured by the maximum number of points in a chain, have a common
finite bound $h$, then $X$ is bi-Esakia representable. Crown cardinalities need
not be uniformly bounded.
\end{corollary}

\begin{proof}
Choose $n$ with $n+1>h$. Since $H_n$ has height $n+1$, it cannot embed into any
crown. Apply Theorem~\ref{thm:crowns}.
\end{proof}

\begin{corollary}[unbounded depth]
Theorem~\ref{thm:crowns} applies to families of crowns with unbounded depth
whenever they uniformly omit one fixed comb. In particular, arbitrary families
of finite chains are covered.
\end{corollary}

\begin{proof}
The fixed poset $H_2$ is branched, so it does not embed into a chain,
independently of the chain length.
\end{proof}

\section{The finite-anchor boundary}

The uniform comb bound is not a decorative assumption in the proof.
Bezhanishvili, Martins, and Moraschini record a countably infinite antichain of
finite co-trees under the surjective bi-p-image relation
\cite[Proposition 4.15]{ref3}. Taking order duals gives an antichain of finite
rooted trees.

If all these rooted trees are placed as crowns above a common root, the
resulting family has no finite set of minimal anchor types generated by the WPO
argument. Thus Theorem~\ref{thm:synchronization} cannot be applied by the
present finite-group reduction.

This is only a limitation of the method. It does not prove that the common-root
tree built from the antichain lacks every possible bi-Esakia topology. A
genuine negative result would require a topology-independent obstruction, or a
proof that any bi-Esakia topology forces a finite quotient comparability absent
from the antichain. Neither statement is established here.

\section{Relation to prior work}

Tasiou's inverse-limit theorem \cite[Theorem 4.2.6]{ref1} supplies the general
categorical background. Her synchronization result for arbitrary copies of one
finite connected poset \cite[Theorem 4.3.4]{ref1}\cite[Theorem 4.6]{ref2} is
the closest explicit construction;
\cite[Theorem 4.3.8]{ref1}\cite[Theorem 4.3]{ref2} instead realizes the finite
stages of a countable inverse system as the zigzag image-finite part of a larger
space. Theorem~\ref{thm:synchronization} differs in three coupled features: the
pieces may be nonisomorphic, only finitely many actual quotient anchors are
required, and all pieces may share a finite downset. The inverse-limit mechanism
itself is not new; the contribution is the stem-exact finite-anchor criterion
and its application to a natural infinite class of bare trees.

Fornasiere and Moraschini \cite{ref4} address Esakia representability of
well-ordered forests. Their result motivates the object class but does not
provide a common bi-Esakia topology. Martins \cite{ref5} proves the BPO theorem
used in Section~\ref{sec:crowns} but does not discuss bi-Esakia topologies on
the common-root amalgams. The antichain in \cite{ref3} explains why the same
finite-anchor argument cannot simply discard the comb restriction.

The present result is therefore a sufficient theorem, not a classification. It
neither proves that the comb condition is necessary nor settles bi-Esakia
representability for arbitrary well-ordered trees or forests.

\section{Open directions}

The most direct continuation is to replace the finite anchor family in
Theorem~\ref{thm:synchronization} by a compact parameter space of anchor types.
Such a construction would have to control clopen up- and down-closures across
infinitely many groups and prevent new inverse-limit points. The known
antichain shows that order-theoretic minimality alone cannot provide this
compactness.

A different route is negative: identify a topology-independent invariant that
prevents a common-root amalgam from being bi-Esakia representable. At present,
the antichain supplies no such invariant.

Other extensions---an infinite common stem, infinite petals, or arbitrary
forest components---each break a separate finiteness step in
Theorem~\ref{thm:synchronization} and should be treated as new problems rather
than formal corollaries.

\section{AI-assisted research and writing disclosure}

This work was developed with AI assistance for literature search, proof
exploration, adversarial checking, computation, and drafting. Carptopus is the
responsible author and independently controls the mathematical claims and
public version. OpenAI Codex is a research and writing tool, not an author.

\begingroup
\small
\raggedright
\setlength{\parskip}{0.25em}
\begin{thebibliography}{9}

\bibitem{ref1}
L. Tasiou,
\emph{Profinite bi-Heyting algebras}, MSc thesis, Korteweg-de Vries Institute
for Mathematics, University of Amsterdam (2021).

\bibitem{ref2}
L. Tasiou,
\emph{Profinite bi-Heyting algebras}, Algebra Universalis 86(3), article 16
(2025).
\url{https://doi.org/10.1007/s00012-025-00892-w}

\bibitem{ref3}
N. Bezhanishvili, M. Martins, and T. Moraschini,
\emph{Bi-intermediate logics of trees and co-trees}, Annals of Pure and Applied
Logic 175 (2024), 103490.
\url{https://doi.org/10.1016/j.apal.2024.103490}

\bibitem{ref4}
D. Fornasiere and T. Moraschini,
\emph{Trees and spectra of Heyting algebras}, Indagationes Mathematicae 37(4)
(2026), 1158--1191.
\url{https://doi.org/10.1016/j.indag.2025.08.002}

\bibitem{ref5}
M. Martins,
\emph{There are only countably many locally tabular bi-intermediate logics of
co-trees}, arXiv:2602.21960v1 (2026).
\url{https://arxiv.org/abs/2602.21960}

\end{thebibliography}
\endgroup

\end{document}
